% TikZ 图：安培力平衡例题示意
% Day 21：磁场与安培力
\begin{tikzpicture}[scale=0.82,>=Latex]
  \foreach \x in {-1.8,-1.1,1.0,1.7} {
    \draw[blue,thick] (\x-0.08,0.67) -- (\x+0.08,0.83);
    \draw[blue,thick] (\x-0.08,0.83) -- (\x+0.08,0.67);
  }
  \foreach \x in {-1.8,-1.1,1.7} {
    \draw[blue,thick] (\x-0.08,-0.78) -- (\x+0.08,-0.62);
    \draw[blue,thick] (\x-0.08,-0.62) -- (\x+0.08,-0.78);
  }
  \node[blue] at (2.35,0.82) {$B$ 向里};
  \draw[very thick,red,->] (-1.45,0) -- (1.45,0) node[right=2pt] {$I$};
  \draw[->,purple,very thick] (0,0) -- (0,1.08) node[above] {$F_A$};
  \draw[->,black,thick] (0,0) -- (0,-0.90) node[pos=0.72,right=4pt] {$mg$};
  \node[font=\small] at (0,-1.45) {悬浮条件：$F_A=mg$};
\end{tikzpicture}
